% titel: Blatt 10, A8: Essigsäure/Acetat-Puffer
% skript: Mortimer Kap. 17 (Puffergleichung)

% Teil I: Werkzeuge, kein Lösungsweg.

Ein sauerer Puffer aus schwacher Säure (Essigsäure) und ihrer konjugierten
Base (Natriumacetat) wird über die Massen angesetzt.

\begin{werkzeugkasten}[Puffer-pH aus Massen]
Aus Masse und Molmasse die Stoffmengen berechnen ($n=m/M$), dann das
Molverhältnis in die Henderson-Hasselbalch-Gleichung einsetzen. Das
Volumen kürzt sich heraus (gleicher V), wichtig sind die Stoffmengen
[Ac$^-$] (aus NaAc) und [HAc].
\[
  pH = pK_s + \lg\frac{n(\mathrm{Ac^-})}{n(\mathrm{HAc})},
  \qquad pK_s(\mathrm{Essigs\text{a}ure})\approx4.75 .
\]
\textit{Skript:} Mortimer Kap. 17 (Puffer).
\end{werkzeugkasten}

\begin{achtung}[Darauf achten]
\begin{itemize}
  \item Molmassen unterscheiden: Essigsäure CH$_3$COOH $\approx60.05$;
        Natriumacetat CH$_3$COONa $\approx82.03$ g/mol.
  \item 50 mg Essigsäure und 10 mg Natriumacetat in mmol umrechnen.
  \item Das Volumen (100 mL) kürzt sich in der Ratio; nur die Mengen zählen.
  \item Ergebnis pH sollte unter $pK_s$ liegen (Säure lag im Überschuss).
\end{itemize}
\end{achtung}
